% Mathématiques complémentaires — Terminale — V3 LaTeX
% Temps d'attente
\section{Objectifs}
\begin{itemize}
\item Modéliser un temps d'attente discret par une loi géométrique.
\item Modéliser un temps d'attente continu par une loi exponentielle.
\item Utiliser la propriété d'absence de mémoire.
\item Calculer et interpréter des probabilités de durée.
\end{itemize}
\section{Repères et enjeux du thème}
Un temps d'attente peut être compté en nombre d'essais ou mesuré continûment. La loi géométrique décrit le rang du premier succès dans une répétition indépendante d'épreuves de Bernoulli. La loi exponentielle modélise certaines durées continues lorsque l'intensité instantanée reste constante. Les deux modèles possèdent une forme d'absence de mémoire.
\section{1. Loi géométrique}
On répète des épreuves de Bernoulli indépendantes de paramètre $p$ jusqu'au premier succès. Si $X$ désigne le rang du premier succès, alors pour $k\geq1$,
\[
P(X=k)=(1-p)^{k-1}p.
\]
Il faut donc obtenir $k-1$ échecs puis un succès. L'espérance vaut
\[
E(X)=\frac1p.
\]
Cette moyenne s'exprime en nombre d'essais.
\section{2. Probabilité de dépasser un rang}
L'événement $X>n$ signifie qu'aucun succès n'a été obtenu au cours des $n$ premiers essais. Ainsi
\[
P(X>n)=(1-p)^n.
\]
On en déduit
\[
P(X\leq n)=1-(1-p)^n.
\]
Ces formules sont utiles pour chercher un nombre minimal d'essais garantissant une probabilité de succès cumulée donnée.
\coursefigure{/images/mathematiques/terminale-mathematiques-complementaires/temps-attente-1.svg}{Distribution géométrique.}{Des bâtons décroissants représentent la probabilité que le premier succès apparaisse aux rangs successifs.}
\section{3. Absence de mémoire discrète}
La loi géométrique vérifie
\[
P(X>m+n\mid X>m)=P(X>n).
\]
Savoir qu'aucun succès n'a eu lieu pendant les $m$ premiers essais ne modifie pas la loi du nombre d'essais supplémentaires à attendre. Cette propriété repose directement sur l'indépendance des épreuves.
Elle ne doit pas être utilisée lorsque la probabilité de succès évolue avec le nombre d'essais.
\section{4. Loi exponentielle}
Une variable aléatoire continue $T$ suit une loi exponentielle de paramètre $\lambda>0$ lorsque sa densité est
\[
f(t)=\lambda e^{-\lambda t},\qquad t\geq0,
\]
et nulle pour $t<0$. La fonction de survie vaut
\[
P(T>t)=e^{-\lambda t}.
\]
La probabilité d'une attente entre $a$ et $b$ s'écrit
\[
P(a\leq T\leq b)=e^{-\lambda a}-e^{-\lambda b}.
\]
\section{5. Espérance et paramètre}
Pour une loi exponentielle,
\[
E(T)=\frac1\lambda.
\]
Le paramètre $\lambda$ représente une intensité et possède l'unité inverse de celle du temps. Si le temps est mesuré en heures, $\lambda$ s'exprime en heure inverse.
Augmenter $\lambda$ raccourcit les temps moyens et accélère la décroissance de la fonction de survie.
\coursefigure{/images/mathematiques/terminale-mathematiques-complementaires/temps-attente-2.svg}{Densité et fonction de survie exponentielles.}{Deux courbes décroissantes illustrent la densité et la probabilité de dépasser un temps donné.}
\section{6. Absence de mémoire continue}
Pour $s\geq0$ et $t\geq0$,
\[
P(T>s+t\mid T>s)=P(T>t).
\]
Le temps déjà attendu n'influence pas la distribution du temps restant dans le modèle. Cette propriété caractérise la loi exponentielle parmi les lois continues usuelles.
Dans une situation réelle, elle suppose que le mécanisme ne vieillit pas et que le risque instantané reste constant. Un composant mécanique soumis à l'usure peut donc ne pas suivre ce modèle.
\section{7. Médiane et quantiles}
La médiane $m$ vérifie $P(T>m)=1/2$. On obtient
\[
e^{-\lambda m}=\frac12,
\qquad
m=\frac{\ln2}{\lambda}.
\]
La médiane diffère de la moyenne $1/\lambda$. Cette différence montre qu'une distribution exponentielle est dissymétrique avec une longue traîne vers les grandes valeurs.
\section{8. Choisir le bon modèle}
La loi géométrique convient lorsque l'on compte des essais distincts jusqu'au premier succès. La loi exponentielle convient lorsque la durée est continue et que l'absence de mémoire est plausible.
Une durée d'attente mesurée à la seconde peut être continue dans le modèle même si l'appareil arrondit l'affichage. Le choix du modèle dépend du mécanisme, pas seulement du format des données.
\section{9. Algorithmique et simulation}
Une simulation géométrique répète des Bernoulli jusqu'au premier succès. Pour une loi exponentielle, un logiciel peut générer directement des durées ou utiliser une transformation adaptée d'une variable uniforme.
Un histogramme de nombreuses durées exponentielles montre une forte concentration près de $0$ et une traîne à droite. La simulation permet de vérifier une moyenne empirique ou une probabilité, sans remplacer la formule théorique.
\section{10. Méthode de résolution}
On commence par identifier si l'attente est discrète ou continue. Pour une loi géométrique, on écrit clairement le nombre d'échecs avant le premier succès. Pour une loi exponentielle, on utilise de préférence la fonction de survie lorsque la question porte sur un dépassement.
Une équation de seuil exponentielle conduit naturellement au logarithme. Les unités doivent être contrôlées et une probabilité finale doit rester comprise entre $0$ et $1$.
\section{11. Lecture critique}
L'absence de mémoire est une hypothèse forte. Elle peut convenir pour certaines arrivées ou certains événements indépendants, mais elle n'est pas une vérité générale sur les durées de vie.
La moyenne seule ne suffit pas à décrire une attente. La médiane, les quantiles et la probabilité de dépasser un seuil apportent des informations complémentaires.
\section{12. Intensité, risque instantané et validation du modèle}
Dans une loi exponentielle, le paramètre $\lambda$ peut se lire comme une intensité constante. Sur un petit intervalle de durée $h$, la probabilité conditionnelle qu'un événement survienne pendant cet intervalle, sachant qu'il n'est pas encore survenu, est approximativement proportionnelle à $\lambda h$. Cette stabilité du risque instantané explique l'absence de mémoire.

Cette interprétation fournit un critère de validation. Si les données montrent que le risque augmente avec l'âge d'un composant ou diminue avec l'expérience d'un utilisateur, une intensité constante devient discutable. La loi exponentielle peut rester une approximation locale, mais il faut éviter de présenter ses propriétés comme exactes pour le phénomène réel.

La fonction de survie permet aussi de comparer deux modèles. Pour des paramètres $\lambda_1<\lambda_2$, on a
\[
e^{-\lambda_1 t}>e^{-\lambda_2 t}\qquad(t>0).
\]
Le premier modèle possède donc une probabilité plus grande de dépasser chaque durée positive et une espérance $1/\lambda_1$ plus élevée. La comparaison peut ainsi être menée sans calculer une multitude de probabilités particulières.

Dans les applications, la cohérence des unités est essentielle. Convertir un paramètre donné par heure en un temps exprimé en minutes demande une conversion avant utilisation. Une erreur d'unité peut produire une valeur numérique plausible mais fausse de plusieurs ordres de grandeur.

\section{Erreurs fréquentes}
\begin{itemize}
\item Confondre rang du premier succès et nombre d'échecs.
\item Oublier que la loi géométrique est discrète et la loi exponentielle continue.
\item Utiliser l'absence de mémoire lorsque le risque évolue avec le temps.
\item Confondre moyenne et médiane.
\end{itemize}
\section{À retenir}
\begin{aretenir}
La loi géométrique modélise une attente discrète jusqu'au premier succès ; la loi exponentielle modélise certaines attentes continues. Toutes deux possèdent une propriété d'absence de mémoire sous des hypothèses précises.
\end{aretenir}
